\documentclass[12pt]{article}
\usepackage{amsmath, amssymb}
\usepackage{fancyhdr}
\usepackage{fullpage}
\usepackage{listings}
\usepackage{enumitem}
\usepackage{graphicx}
\usepackage{courier}
\usepackage{float}
\usepackage{color}


\title{CPSC 250 - Programming for Data Manipulation \\ Final Exam – Solution Set}
\date{}
\begin{document}
\maketitle

\section*{Part I: Multiple Choice (30 points)}

\textbf{Correct Answers:}

\begin{enumerate}[label=\arabic*.]
\item \textbf{B} – Lists are mutable, and \texttt{b} and \texttt{a} point to the same object.
\item \textbf{B} – \texttt{dropna()} removes missing data rows or columns.
\item \textbf{B} – \texttt{\_\_str\_\_} defines the string representation.
\item \textbf{C} – Inheritance promotes code reuse.
\item \textbf{C} – The base case defines when recursion stops.
\item \textbf{D} – Lists are mutable; tuples are immutable.
\item \textbf{B} – \texttt{curve\_fit} finds parameters that best fit a function.
\item \textbf{C} – Balanced BSTs allow \texttt{O(log n)} searches.
\item \textbf{A} – \texttt{\{\}} defines an empty dictionary.
\item \textbf{A} – \texttt{\_\_eq\_\_} customizes \texttt{==} behavior.
\item \textbf{D} – \texttt{read\_csv()} returns a DataFrame.
\item \textbf{C} – Operator overloading allows operators to work with user types.
\item \textbf{D} – \texttt{.mean()} returns the average.
\item \textbf{B} – Polymorphism means shared interface, different behavior.
\item \textbf{C} – Inheritance uses syntax \texttt{class A(B):}
\end{enumerate}

\newpage
\section*{Part II: Error Identification (15 points)}

\textbf{Original Code:}
\begin{lstlisting}[language=Python]
def reverse_list(lst):
    if len(lst) = 0:
        return []
    else:
        return reverse_list(lst[1:]) + lst[0]
\end{lstlisting}

\textbf{Corrected Code:}
\begin{lstlisting}[language=Python]
def reverse_list(lst):
    if len(lst) == 0:
        return []
    else:
        return reverse_list(lst[1:]) + [lst[0]]
\end{lstlisting}

\textbf{Explanation:}
\begin{itemize}
\item The condition should use \texttt{==}, not \texttt{=}.
\item You must concatenate lists, not a list and a single value. So use \texttt{[lst[0]]}.
\end{itemize}

\newpage
\section*{Part III: Code Writing (30 points)}

\textbf{Q1. Palindrome Checker (6 points)}
\begin{lstlisting}[language=Python]
def is_palindrome(s):
    return s == s[::-1]
\end{lstlisting}

---

\textbf{Q2. Student class with \_\_str\_\_ and \_\_lt\_\_ (8 points)}
\begin{lstlisting}[language=Python]
class Student:
    def __init__(self, name, gpa):
        self.name = name
        self.gpa = gpa

    def __str__(self):
        return f"{self.name}, GPA: {self.gpa}"

    def __lt__(self, other):
        return self.gpa < other.gpa
\end{lstlisting}

---

\textbf{Q3. Load and filter DataFrame (8 points)}
\begin{lstlisting}[language=Python]
import pandas as pd

def load_and_filter(filename):
    df = pd.read_csv(filename)
    return df[df['score'] > 80]
\end{lstlisting}

---

\textbf{Q4. Animal classes with polymorphism (8 points)}
\begin{lstlisting}[language=Python]
class Animal:
    def make_sound(self):
        return "..."

class Dog(Animal):
    def make_sound(self):
        return "woof"

class Cat(Animal):
    def make_sound(self):
        return "meow"

# Demonstrate polymorphism
animals = [Dog(), Cat(), Dog()]
for animal in animals:
    print(animal.make_sound())
\end{lstlisting}

\newpage
\section*{Part IV: Code Comprehension (25 points)}

\textbf{Original Code:}
\begin{lstlisting}[language=Python]
import numpy as np
from scipy.optimize import curve_fit

def model(x, a, b):
    return a * np.exp(b * x)

def fit_data(x, y):
    popt, pcov = curve_fit(model, x, y)
    print("a =", popt[0])
    print("b =", popt[1])
\end{lstlisting}

\textbf{Line-by-line Comments:}
\begin{itemize}
\item \texttt{import numpy as np} Load NumPy for numeric operations.
\item \texttt{from scipy.optimize import curvefit} Import curve fitting function.
\item \texttt{def model(x, a, b):} Define a model function with two parameters.
\item \texttt{return a * np.exp(b * x)} Return exponential function output.
\item \texttt{def fitdata(x, y):} Define function that fits data to the model.
\item \texttt{popt, pcov = curvefit(...)} Fit model to data, returning parameters and covariances.
\item \texttt{print("a =", popt[0])} Print the best-fit value for \texttt{a}.
\item \texttt{print("b =", popt[1])} Print the best-fit value for \texttt{b}.
\end{itemize}

\textbf{What do \texttt{popt} and \texttt{pcov} represent?}
\begin{itemize}
\item \texttt{popt}: A list of optimal values for the parameters (\texttt{a}, \texttt{b}).
\item \texttt{pcov}: The estimated covariance matrix of the parameter estimates.
\end{itemize}

\end{document}
